%%%%%%%%%%%%
%
% $Autor: Wings $
% $Datum: 2019-03-05 08:03:15Z $
% $Pfad: TemplateSensor $
% $Version: 4250 $
% !TeX spellcheck = en_GB/de_DE
% !TeX encoding = utf8
% !TeX root = filename 
% !TeX TXS-program:bibliography = txs:///biber
%
%%%%%%%%%%%%

% Structure
\chapter{MKR-Erweiterungsplatine}

Dieses Kapitel beschreibt die MKR-Erweiterungsplatine als Schnittstelle zur einfachen Anbindung von Sensoren, Aktoren und Grove-Modulen an ein Arduino-MKR-Board. Es behandelt deren grundlegenden Aufbau, die integrierte Spannungsversorgung über den Buck-Wandler sowie die verfügbaren Ein- und Ausgänge. Zudem werden die wichtigsten technischen Spezifikationen und die Nutzung im praktischen Einsatz erläutert.


\section{General}
	Die MKR-Erweiterungsplatine wurde entwickelt, um die Nutzung von Arduino MKR-Boards zu vereinfachen. Sie bietet standardisierte Anschlüsse für Grove-Sensoren sowie andere Module und erleichtert so den Aufbau von IoT-, Automatisierungs- oder Steuerungsprojekten. Dank der vorkonfigurierten Schnittstellen entfällt aufwendiges Verdrahten, sodass sich Prototypen schneller umsetzen lassen.
	
	Ein wesentliches Merkmal der Erweiterungsplatine ist ein Buck-Wandler, der mit einer externen Spannung von 7–16V über den VIN-Anschluss der Schraubklemme versorgt werden kann. Dieser Wandler liefert die 5V- und 3,3V-Spannungen, die ebenfalls an den Schraubklemmen verfügbar sind. Allerdingswerden sowohl alle digitale Ein- und Ausgänge, als auch Analoge Eingänge mit 5V versorgt.

	\begin{center}
		%todo tikz picture with strip red - red - black - brown 1\%
		% or red - red - black - gold 5\%
		\includegraphics[width=0.8\textwidth]{MKRWAN1310/MKRConnectorCarrier}
		\captionof{figure}{MKR Connector Carrier}
	\end{center}

%\section{Specific Sensor}

\section{Specification}

	\begin{table}[H]
		\centering
		\begin{tabular}{|l|l|}
			\hline
			\textbf{Eigenschaft} & \textbf{Wert} \\
			\hline
			Eingangsspannung & 7 V - 16 V \\
			\hline
			Analoge Eingänge & 7 \\
			\hline
			Digitale I/Os & 7 \\
			\hline
			Schnittstellen & UART, I2C \\
			\hline
			Stromverbrauch & 5 mA \\
			\hline
			Gewicht & 25 g \\
			\hline
			Produktcode & ASX00007 \\
			\hline
		\end{tabular}
		\caption{Technische Spezifikationen des MKR Connector Carrier}
	\end{table}
	
\section{Bibliothek}
	Das Shield erfordert keine zusätzlichen Bibliotheken. Es funktioniert Plug \& Play, d. h., das Arduino muss in das Shield gesteckt werden.

%\subsection{Description}
\section{Installation}
	Der Arduino soll in das Shield gesteckt werden. Auf dem Shield befinden sich bei den Anschlüssen Beschreibungen, welcher Anschluss mit welchem Arduino-Pin verbunden ist. Die entsprechenden Sensoren sollen folglich an den gewünschten Steckdosen angeschlossen werden. Die jeweiligen Pins sollen dann im Programm initialisiert und definiert werden.

\subsection{Functions}

%\subsection{Example - Manual}
%\subsection{Example}
%\subsection{Example - Code}
%\subsection{Example - Files}



%\section{Calibration}
%\section{Simple Code}
%\section{Simple Application}
%\section{Tests}
%\subsection{Simple Function Test}
%\subsection{Test all Functions}
%\section{Simple Application}
%\section{Further Readings}

